\section{Derive the Aguilar--Balslev--Combes mechanism step by step}
\label{app:abc-sketch}
% BEGIN CURATED SUBJECT INDEX
\index{operator!adjoint and self-adjoint}
\index{spectrum!point}
\index{resolvent}
\index{resolvent set}
\index{resonance!residue}
\index{scattering!Coulomb}
\index{Aguilar--Balslev--Combes theorem}
% END CURATED SUBJECT INDEX

This chapter supplies the functional-analytic chain behind the statement in
Section~\ref{subsec:abc}.  We first state one precise one-threshold version;
the following steps prove its mechanism.  Operator-domain and quadratic-form
variants enlarge the admissible class of potentials
\cite{Aguilar:1971ve,Balslev:1971vb,Simon1972BalslevCombes,Simon:73}.

\begin{theorem}[Aguilar--Balslev--Combes, type-A one-threshold form]
Let $H_0=-\hbar^2\Delta/(2\mu)$ on $L^2(\real^d)$ with domain $H^2(\real^d)$,
and let $H=H_0+V$ be self-adjoint.  Suppose there is $\theta_0>0$ such that
for $|\im s|<\theta_0$ the dilation
$V_s=U(s)VU(-s)$ is analytic as an operator from $D(H_0)$ with its graph norm
to $L^2$, and $V_s(H_0+1)^{-1}$ is compact locally uniformly in $s$.  Assume
also that $H(s)=\rme^{-2s}H_0+V_s$ is a closed type-A family on the common
domain $D(H_0)$.

For $0<\theta<\theta_0$, put $H_\theta=H(\rmi\theta)$.  Then
\begin{enumerate}[label=(\roman*)]
  \item $\essspectrum(H_\theta)
  =\rme^{-2\rmi\theta}[0,\infty)$;
  \item the remaining spectrum away from this ray consists of isolated
  eigenvalues of finite algebraic multiplicity;
  \item matrix elements $\dualpair{\phi}{(z-H)^{-1}\psi}$ between
  dilation-analytic vectors continue meromorphically from the physical
  resolvent set into the wedge swept by the ray, and their nonzero-residue
  poles coincide with the discrete eigenvalues of $H_\theta$ there; and
  \item these eigenvalues and their algebraic multiplicities are independent
  of $\theta$ while they remain in a common uncovered wedge.  Negative
  isolated eigenvalues of $H$ remain unchanged.
\end{enumerate}
\label{thm:abc-type-A}
\end{theorem}

The assumptions are deliberately visible.  Coulomb tails, a potential that
is not analytic in the swept coordinate sector, and threshold accumulation
require modified theorems rather than an invocation of the displayed result.

\subsection*{Step 1: real dilations and analytic vectors}
% BEGIN CURATED SUBJECT INDEX
\index{resolvent}
% END CURATED SUBJECT INDEX

For real $s$, the map
\begin{equation}
  [U(s)\psi](\bm x)=\rme^{ds/2}\psi(\rme^s\bm x)
  \label{eq:abc-app-real-dilation}
\end{equation}
is unitary on $L^2(\real^d)$ and satisfies
\begin{equation}
  U(s)H_0U(-s)=\rme^{-2s}H_0 .
  \label{eq:abc-app-free-scaling}
\end{equation}
The continuation $s\to\rmi\theta$ is not a bounded unitary operation on all
of $L^2$.  One therefore chooses a dense set $\mathcal A$ of vectors for
which $s\mapsto U(s)\phi$ is analytic in a strip.  Gaussian functions and
finite linear combinations of oscillator eigenfunctions provide standard
examples.  The vector set is used to continue matrix elements; the operator
$H(\rmi\theta)$ is defined independently as a closed analytic-family member.

Assume that $V_s=U(s)VU(-s)$ continues analytically and that
\begin{equation}
  V_s(H_0+1)^{-1}
  \label{eq:abc-app-relative-compact}
\end{equation}
is compact, locally uniformly in the analytic strip.  A common-domain
construction gives a type-A family
\begin{equation}
  H(s)=\rme^{-2s}H_0+V_s,
  \qquad D(H(s))=D(H_0).
  \label{eq:abc-app-type-A}
\end{equation}
Singular interactions are often better handled by analytic quadratic forms,
which give a type-B family.  Subsequent resolvent arguments are structurally
the same.

\subsection*{Step 2: the rotated free spectrum}
% BEGIN CURATED SUBJECT INDEX
\index{operator!adjoint and self-adjoint}
\index{resolvent}
\index{compact operator}
\index{Fredholm theory}
% END CURATED SUBJECT INDEX

For $s=\rmi\theta$,
\begin{equation}
  H_{0,\theta}=\rme^{-2\rmi\theta}H_0 .
  \label{eq:abc-app-free-Htheta}
\end{equation}
The spectrum of $H_0$ is $[0,\infty)$, so multiplication by the complex
number $\rme^{-2\rmi\theta}$ gives
\begin{equation}
  \spectrum(H_{0,\theta})
  =\rme^{-2\rmi\theta}[0,\infty) .
  \label{eq:abc-app-free-spectrum}
\end{equation}
This is the elementary source of the angle $2\theta$: the coordinate scales
as $\rme^{\rmi\theta}$, momentum as $\rme^{-\rmi\theta}$, and kinetic energy
as the square of momentum.

For $z$ off the free ray, the second resolvent identity in the convention
$\Resolvent(z)=(z-H)^{-1}$ gives
\begin{align}
  &(z-H_\theta)^{-1}-(z-H_{0,\theta})^{-1}
  \notag\\
  &\qquad=(z-H_\theta)^{-1}V_\theta
  (z-H_{0,\theta})^{-1} .
  \label{eq:abc-app-resolvent-difference}
\end{align}
Under the relative-compactness hypothesis, the right-hand side is compact
whenever both resolvents exist.  With the Fredholm definition of essential
spectrum appropriate to a non-self-adjoint operator, compact resolvent
perturbations do not change that spectrum.  Hence
\begin{equation}
  \essspectrum(H_\theta)
  =\rme^{-2\rmi\theta}[0,\infty) .
  \label{eq:abc-app-essential-spectrum}
\end{equation}

\subsection*{Step 3: analytic Fredholm theory}
% BEGIN CURATED SUBJECT INDEX
\index{spectrum!point}
\index{resolvent}
\index{Fredholm theory}
\index{Jordan chain}
% END CURATED SUBJECT INDEX

Away from the rotated ray, factor
\begin{equation}
  z-H_\theta
  =\left[\identity-V_\theta(z-H_{0,\theta})^{-1}\right]
  (z-H_{0,\theta}) .
  \label{eq:abc-app-fredholm-factor}
\end{equation}
The operator
\begin{equation}
  K_\theta(z)=V_\theta(z-H_{0,\theta})^{-1}
  \label{eq:abc-app-K}
\end{equation}
is compact and analytic in $z$.  For sufficiently negative real $z$, or at
another point controlled by a resolvent estimate,
$\identity-K_\theta(z)$ is invertible.  The analytic Fredholm theorem then
states that
\begin{equation}
  [\identity-K_\theta(z)]^{-1}
  \label{eq:abc-app-fredholm-inverse}
\end{equation}
is meromorphic away from the rotated ray.  Its singularities are isolated,
and the principal parts have finite rank.  By
Eq.~\eqref{eq:abc-app-fredholm-factor}, these singularities are precisely the
discrete eigenvalues of $H_\theta$, counted with algebraic multiplicity.

If $\gamma$ is a small contour around one such eigenvalue $z_n$, the Riesz
projection is
\begin{equation}
  P_n(\theta)
  =\frac{1}{2\pi\rmi}\oint_\gamma
  (z-H_\theta)^{-1}\rmd z .
  \label{eq:abc-app-riesz-projection}
\end{equation}
Its rank is the algebraic multiplicity.  This formulation includes
nontrivial Jordan blocks; diagonalizability is not assumed.

\subsection*{Step 4: continuation of physical resolvent matrix elements}
% BEGIN CURATED SUBJECT INDEX
\index{Hilbert space}
\index{spectrum!point}
\index{resolvent}
\index{GNS construction}
\index{resonance!residue}
% END CURATED SUBJECT INDEX

Let $H(s)=U(s)HU(-s)$ for real $s$.  Unitarity gives, for
$\phi,\psi\in\mathcal A$,
\begin{equation}
  \bra{\phi}(z-H)^{-1}\ket{\psi}
  =\bra{U(s)\phi}(z-H(s))^{-1}U(s)\ket{\psi} .
  \label{eq:abc-real-identity}
\end{equation}
To continue this identity holomorphically, one must remember that the inner
product is antilinear in its first argument.  The analytic expression is
\begin{equation}
  \Dualpair{U(\overline{s})\phi}
  {(z-H(s))^{-1}U(s)\psi} .
  \label{eq:abc-sesquilinear-analytic-expression}
\end{equation}
For $z$ initially above the physical spectrum, it extends jointly in $(s,z)$
to an overlapping complex domain and agrees with the physical matrix element
for real $s$.  The identity theorem permits continuation to
$s=\rmi\theta$:
\begin{equation}
  \bra{\phi}(z-H)^{-1}\vsub{\ket{\psi}}{continued}
  =\bra{U(-\rmi\theta)\phi}
  (z-H_\theta)^{-1}U(\rmi\theta)\ket{\psi} .
  \label{eq:abc-app-continued-matrix-element}
\end{equation}
The left-hand side is notation for a continued scalar matrix element, not a
bounded resolvent on the original Hilbert space.  The right-hand side is
meromorphic by Step 3.  Therefore a discrete eigenvalue of $H_\theta$ in the
uncovered wedge is a pole of the continued physical matrix element.

The opposite signs in the two analytic vectors are forced by sesquilinearity;
using $U(\rmi\theta)$ in both slots would continue a different expression.

Conversely, a pole that couples to analytic vectors appears as an eigenvalue
once the rotated ray lies below it.  A particular pair $\phi,\psi$ can have a
zero residue and fail to see a pole, but the pole is detected by some vectors
in a dense analytic set.

\subsection*{Step 5: angle independence and bound states}
% BEGIN CURATED SUBJECT INDEX
\index{operator!adjoint and self-adjoint}
\index{spectrum!point}
\index{Riesz projection}
% END CURATED SUBJECT INDEX

Take two admissible angles $\theta_1<\theta_2$ and a domain of the energy
plane uncovered by both rotated rays.  Equation
\eqref{eq:abc-app-continued-matrix-element} represents the same analytic
continuation for either angle.  Uniqueness of analytic continuation therefore
prevents an isolated pole from moving with $\theta$.  Applying the argument
to matrix elements of the Riesz projection shows that its rank, hence the
algebraic multiplicity, is also constant until the pole meets the essential
spectrum or exits the analyticity domain.

A negative-energy bound state is already an isolated eigenvalue of the
undeformed self-adjoint Hamiltonian.  For real $s$ it is unchanged because
$H(s)$ is unitarily equivalent to $H$.  Analyticity in $s$ then continues this
constant eigenvalue to complex $s$.  This is the mathematical origin of the
stationary real points in a complex-scaling plot.

\subsection*{Step 6: several thresholds and the limits of the theorem}
% BEGIN CURATED SUBJECT INDEX
\index{spectrum!essential}
\index{compact operator}
\index{Fredholm theory}
\index{complex scaling}
\index{complex scaling!exterior}
\index{tortoise coordinate}
\index{horizon boundary condition}
% END CURATED SUBJECT INDEX

For an $N$-body Hamiltonian, cluster decompositions produce thresholds
$E_a$.  Relative motion in an open cluster channel supplies a kinetic
continuum, so the corresponding essential spectrum rotates about its own
threshold:
\begin{equation}
  \essspectrum(H_\theta)
  =\bigcup_a\left[E_a+\rme^{-2\rmi\theta}[0,\infty)\right]
  \label{eq:abc-app-many-thresholds}
\end{equation}
under the hypotheses of the many-body theorem.  A pole must be classified on
the multi-sheeted surface generated by these thresholds.

The argument also displays its failure modes.  They include a potential with
no analytic continuation in the swept sector, loss of relative compactness,
a deformation crossing a singularity, a long-range remainder not separated
from the asymptotic Hamiltonian, and a threshold at which the Fredholm sheet
changes.  Exterior complex scaling relaxes interior analyticity by deforming
only outside a finite radius, but it still requires control of the asymptotic
operator.  For black-hole equations, logarithmic tortoise coordinates,
multiple horizons, inverse-power tails, and frequency-dependent or coupled
potentials must be checked individually.  Successful numerical rotation is
strong evidence, but it is not by itself a proof that all ABC hypotheses hold.
